\documentclass[../../main.tex]{subfiles}
\graphicspath{{\subfix{../../image/}}}

\begin{document}

\section*{习题一}\label{section:ch01-习题一}

\begin{exercise}\label{exercise:ch01-01}
在某城市中，共发行三种报纸 $A,B,C$。在这城市的居民中，订阅 $A$ 的占 $45\%$，订阅 $B$ 的占 $35\%$，订阅 $C$ 的占 $30\%$，同时订阅 $A$ 及 $B$ 的占 $10\%$，同时订阅 $A$ 及 $C$ 的占 $8\%$，同时订阅 $B$ 及 $C$ 的占 $5\%$，同时订阅 $A,B,C$ 的占 $3\%$，试求下列百分率：（1）只订阅 $A$ 的；（2）只订阅 $A$ 及 $B$ 的；（3）只订阅一种报纸的；（4）正好订阅两种报纸的；（5）至少订阅一种报纸的；（6）不订阅任何报纸的。
\end{exercise}

\begin{exercise}\label{exercise:ch01-02}
若 $A,B,C$ 是随机事件，说明下列关系式的概率意义：（1）$ABC=A$；（2）$A\cup B\cup C=A$；（3）$AB\subset C$；（4）$A\subset\overline{BC}$。
\end{exercise}

\begin{exercise}\label{exercise:ch01-03}
在某班学生中任选一个同学，以事件 $A$ 表示选到的是男同学，事件 $B$ 表示选到的人不喜欢唱歌，事件 $C$ 表示选到的人是运动员。（1）表述 $AB\overline{C}$ 及 $A\overline{B}C$；（2）什么条件下成立 $ABC=A$；（3）何时成立 $\overline{C}\subset B$；（4）何时同时成立 $A=B$ 及 $\overline{A}=C$。
\end{exercise}

\begin{exercise}\label{exercise:ch01-04}
试证：$\bigcup_{i=1}^{n}A_i=A_1+\overline{A}_1A_2+\overline{A}_1\overline{A}_2A_3+\cdots+\overline{A}_1\overline{A}_2\overline{A}_3\cdots\overline{A}_{n-1}A_n$，并对 $n=4$，画出文图。
\end{exercise}

\begin{exercise}\label{exercise:ch01-05}
为``剪刀·石头·布''游戏造一个样本空间，定义有关事件，并考虑如何给定概率。
\end{exercise}

\begin{exercise}\label{exercise:ch01-06}
若 $A,B,C,D$ 是四个事件，试用这四个事件表示下列各事件：（1）这四个事件至少发生一个；（2）$A,B$ 都发生而 $C,D$ 都不发生；（3）这四个事件恰好发生两个；（4）这四个事件都不发生；（5）这四个事件中至多发生一个。
\end{exercise}

\begin{exercise}\label{exercise:ch01-07}
从 $0,1,2,\cdots,9$ 中随机地取出 $5$ 个数（可重复），以 $E_i$ 记某些数正好出现 $i$ 次这一事件（例如 $52353$，既属于 $E_1$，也属于 $E_2$ 及 $E_0$），试用文图表示 $E_0,E_1,\cdots,E_6$ 的关系。
\end{exercise}

\begin{exercise}\label{exercise:ch01-08}
证明下列等式：

（1）$\displaystyle\binom{n}{1}+2\binom{n}{2}+3\binom{n}{3}+\cdots+n\binom{n}{n}=n2^{n-1}$

（2）$\displaystyle\binom{n}{1}-2\binom{n}{2}+3\binom{n}{3}+\cdots+(-1)^{n-1}n\binom{n}{n}=0$

（3）$\displaystyle\sum_{k=0}^{a-r}\binom{a}{k+r}\binom{b}{k}=\binom{a+b}{a-r}$
\end{exercise}

\begin{exercise}\label{exercise:ch01-09}
一部五卷的文集，按任意次序放到书架上去，试求下列概率：（1）第一卷出现在旁边；（2）第一卷及第五卷出现在旁边；（3）第一卷或第五卷出现在旁边；（4）第一卷及第五卷都不出现在旁边；（5）第三卷正好在正中。
\end{exercise}

\begin{exercise}\label{exercise:ch01-10}
甲袋中有 $3$ 只白球，$7$ 只红球，$15$ 只黑球，乙袋中有 $10$ 只白球，$6$ 只红球，$9$ 只黑球，现从两袋中各取一球，求两球颜色相同的概率。
\end{exercise}

\begin{exercise}\label{exercise:ch01-11}
袋中有 $n$ 只球，记有号码 $1,2,\cdots,n$，求下列事件的概率：（1）任意取出 $2$ 球，号码为 $1,2$；（2）任意取出 $3$ 球，没有号码 $1$；（3）任意取出 $5$ 球，号码 $1,2,3$ 中至少出现一个。
\end{exercise}

\begin{exercise}\label{exercise:ch01-12}
袋中装有 $1,2,\cdots,N$ 号的球各一只，采用（1）有放回；（2）不放回方式摸球，试求在第 $k$ 次摸球时首次摸到 $1$ 号球的概率。
\end{exercise}

\begin{exercise}\label{exercise:ch01-13}
从 $6$ 双不同的手套中任取 $4$ 只，问其中恰有一双配对的概率是多少？
\end{exercise}

\begin{exercise}\label{exercise:ch01-14}
从 $n$ 双不同的鞋子中任取 $2r$（$2r<n$）只，求下列事件发生的概率：（1）没有成对的鞋子；（2）只有一对鞋子；（3）恰有两对鞋子；（4）有 $r$ 对鞋子。
\end{exercise}

\begin{exercise}\label{exercise:ch01-15}
$m$ 个男孩和 $n$ 个女孩（$n\leqslant m$）随机地沿着圆桌坐下，试求任意两个女孩都不相邻的概率。
\end{exercise}

\begin{exercise}\label{exercise:ch01-16}
（分赌注问题）甲、乙二人各出赌注 $a$，约定谁先胜三局则赢得全部赌注，现已赌三局，甲二胜一负，这时因故中止赌博，若二人赌技相同，问应如何分配赌注，才算公平合理？
\end{exercise}

\begin{exercise}\label{exercise:ch01-17}
从 $52$ 张扑克牌中任意取出 $13$ 张，求：（1）有 $5$ 张黑桃，$3$ 张红心，$3$ 张方块，$2$ 张草花的概率；（2）牌型分布为 $7$-$3$-$2$-$1$（最长花色有 $7$ 张，最短花色有 $1$ 张，其余二花色分别有 $3$ 张及 $2$ 张）的概率。
\end{exercise}

\begin{exercise}\label{exercise:ch01-18}
桥牌游戏中（四人各从 $52$ 张纸牌中分得 $13$ 张），求 $4$ 张 A 集中在一个人手中的概率。
\end{exercise}

\begin{exercise}\label{exercise:ch01-19}
在扑克牌游戏中（从 $52$ 张牌中任取 $5$ 张），求下列事件的概率：（1）以 A 打头的同花顺次五张牌；（2）其他同花顺次五张牌；（3）有四张牌同点数；（4）三张同点数且另两张也同点数；（5）五张同花；（6）异花顺次五张牌；（7）三张同点数，另外两张不同点数；（8）五张中有两对；（9）五张中有一对；（10）其他情况。
\end{exercise}

\begin{exercise}\label{exercise:ch01-20}
从装有号码 $1,2,\cdots,N$ 的球的箱子中有放回地摸了 $n$ 次球，依次记下其号码，试求这些号码按严格上升次序排列的概率。
\end{exercise}

\begin{exercise}\label{exercise:ch01-21}
在上题中这些号码按上升（不一定严格）次序排列的概率。
\end{exercise}

\begin{exercise}\label{exercise:ch01-22}
任意从数列 $1,2,\cdots,N$ 中不放回地取出 $n$ 个数并按大小排列成：$x_1<x_2<\cdots<x_m<\cdots<x_n$，试求 $x_m=M$ 的概率。这里 $1\leqslant M\leqslant N$。
\end{exercise}

\begin{exercise}\label{exercise:ch01-23}
上题中，若采用有放回取数，这时 $x_1\leqslant x_2\leqslant\cdots\leqslant x_m\leqslant\cdots\leqslant x_n$，试求 $x_m=M$ 的概率。
\end{exercise}

\begin{exercise}\label{exercise:ch01-24}
从 $(0,1)$ 中随机地取两个数，求下列概率：（1）两数之和小于 $1.2$；（2）两数之积小于 $\frac{1}{4}$；（3）以上两个要求同时满足。
\end{exercise}

\begin{exercise}\label{exercise:ch01-25}
从 $(0,1)$ 中随机地取二数 $b$ 及 $c$，试求方程 $x^2+bx+c=0$ 有实根的概率。
\end{exercise}

\begin{exercise}\label{exercise:ch01-26}
在一张打上方格的纸上投一枚直径为 $1$ 的硬币，方格要多小才能使硬币与线不相交的概率小于 $1\%$。
\end{exercise}

\begin{exercise}\label{exercise:ch01-27}
某码头只能容纳一只船，现预知某日将独立来到两只船，且在 $24$ 小时内各时刻来到的可能性都相等，如果它们需要停靠的时间分别为 $3$ 小时及 $4$ 小时，试求有一船要在江中等待的概率。
\end{exercise}

\begin{exercise}\label{exercise:ch01-28}
两人约定于 $7$ 点到 $8$ 点在某地会面，试求一人要等另一人半小时以上的概率。
\end{exercise}

\begin{exercise}\label{exercise:ch01-29}
在线段 $AB$ 中随机地取两个点把线段截为三段，求这三段可以构成一个三角形的概率（三线段能构成三角形的充要条件是任意二边之和大于第三边）。
\end{exercise}

\begin{exercise}\label{exercise:ch01-30}
在线段 $[0,1]$ 上任意投三个点，问由 $0$ 至三点的三线段能构成三角形与不能构成三角形这两个事件中哪一个事件的概率大。
\end{exercise}

\begin{exercise}\label{exercise:ch01-31}
从一只装有 $100$ 只灯泡的箱子中任抽 $5$ 只灯泡，发现有 $2$ 只是次品，你对此批灯泡的次品数作何估计？（这种抽查当然用不放回方式。比较用最大似然估计法所得结果与用频率估计概率法的结果是否相同。）
\end{exercise}

\begin{exercise}\label{exercise:ch01-32}
利用概率论的想法证明下列恒等式：
\begin{align*}
1+\frac{A-a}{A-1}+\frac{(A-a)(A-a-1)}{(A-1)(A-2)}+\cdots+\frac{(A-a)\cdots2\cdot1}{(A-1)\cdots(a+1)a}=\frac{A}{a}
\end{align*}
其中 $A,a$ 都是正整数，且 $A>a$。
\end{exercise}

\begin{exercise}\label{exercise:ch01-33}
设 $A_1,A_2,\cdots,A_n$ 是随机事件，试用归纳法证明下列公式
\begin{align*}
P(A_1\cup A_2\cup\cdots\cup A_n)&=\sum_{i=1}^{n}P(A_i)-\sum_{1\leqslant i<j\leqslant n}P(A_iA_j)\\
&\quad+\sum_{1\leqslant i<j<k\leqslant n}P(A_iA_jA_k)-\cdots+(-1)^{n-1}P(A_1A_2\cdots A_n)
\end{align*}
\end{exercise}

\begin{exercise}\label{exercise:ch01-34}
某班有 $N$ 个士兵，每人各有一支枪，这些枪外形完全一样，在一次夜间紧急集合中，若每人随机地取走一支枪，问至少有一个人拿到自己的枪的概率。
\end{exercise}

\begin{exercise}\label{exercise:ch01-35}
在上题中求恰好有 $k$（$0\leqslant k\leqslant N$）个人拿到自己的枪的概率。
\end{exercise}

\begin{exercise}\label{exercise:ch01-36}
从一副扑克牌中有放回地一张张抽取，求在抽取第 $6$ 张时得到全部 $4$ 种花色的概率。
\end{exercise}

\begin{exercise}\label{exercise:ch01-37}
考试时共有 $N$ 张考签，$n$ 个学生参加考试（$n\geqslant N$），被抽过的考签立刻放回，求在考试结束之后，至少有一张考签没有被抽到的概率。
\end{exercise}

\begin{exercise}\label{exercise:ch01-38}
（赠券收集）食品厂把印有水浒 $108$ 将之一的画卡作为赠券装入某种儿童食品袋中，每袋一卡，试求购买 $n$ 袋这种食品而能收齐全套画卡的概率。
\end{exercise}

\begin{exercise}\label{exercise:ch01-39}
用概率论想法求 $N$ 阶行列式的展开式中包含主对角线元素的项数。
\end{exercise}

\begin{exercise}\label{exercise:ch01-40}
有 $w$ 个白球与 $b$ 个黑球任你放入两个袋子中，让你的朋友随机抽一袋并从中摸出一只球，你将如何做以使你的朋友摸得黑球的概率最大。
\end{exercise}

\begin{exercise}\label{exercise:ch01-41}
甲，乙，丙三人按下面规则进行比赛，第一局由甲，乙参加而丙轮空，由第一局的优胜者与丙进行第二局比赛，而失败者则轮空，比赛用这种方式一直进行到其中一个人连胜两局为止，连胜两局者成为整场比赛的优胜者，若甲，乙，丙胜每局的概率各为 $\frac{1}{2}$，问甲，乙，丙成为整场比赛优胜者的概率各是多少？
\end{exercise}

\begin{exercise}\label{exercise:ch01-42}
父，母，子三人举行比赛，每局总有一人胜一人负（没有和局），每局的优胜者就与未参加此局的人再进行比赛，如果某人首先胜了两局，则他就是整个比赛的优胜者，由父决定第一局由哪两人参加，其中儿子实力最强，所以父为了使自己得胜的概率达到最大，就决定第一局由他与妻子先比赛，试证父的决策为最优策略（任何一对选手中一人胜对方的概率在整个比赛中是不变的）。
\end{exercise}

\begin{exercise}\label{exercise:ch01-43}
给定 $p=P(A)$，$q=P(B)$，$r=P(A\cup B)$，求 $P(A\overline{B})$ 及 $P(\overline{A}\ \overline{B})$。
\end{exercise}

\begin{exercise}\label{exercise:ch01-44}
设 $p_1,p_2,p_{12}$ 是给定的实数，试证存在两个事件 $A_1$ 及 $A_2$ 使得 $P(A_1)=p_1$，$P(A_2)=p_2$，$P(A_1A_2)=p_{12}$ 的充要条件是下列四个不等式同时成立：
\begin{align*}
p_{12}\geqslant 0,\quad p_1-p_{12}\geqslant 0,\quad p_2-p_{12}\geqslant 0,\quad 1-p_1-p_2+p_{12}\geqslant 0
\end{align*}
\end{exercise}

\begin{exercise}\label{exercise:ch01-45}
证明：$\left|P(AB)-P(A)P(B)\right|\leqslant\dfrac{1}{4}$，并讨论等号成立的条件。
\end{exercise}

\begin{exercise}\label{exercise:ch01-46}
求包含事件 $A,B$ 的最小 $\sigma$ 域。
\end{exercise}

\begin{exercise}\label{exercise:ch01-47}
证明：（1）$\Omega$ 的一切子集组成的集类是一个 $\sigma$ 域；（2）$\sigma$ 域之交仍为 $\sigma$ 域。
\end{exercise}

\begin{exercise}\label{exercise:ch01-48}
证明：包含一切形为 $(-\infty,x)$ 的区间的最小 $\sigma$ 域是一维 Borel $\sigma$ 域。
\end{exercise}

\begin{exercise}\label{exercise:ch01-49}
（1）设 $Q$ 是定义在 $\sigma$ 域上的非负广义实值函数（即可以取有限或无限值的函数），如果它具有可列可加性，并且 $Q(\varnothing)=0$，则称 $Q$ 为测度，试说明测度概念是算术中计数概念及几何中长度、面积、体积等概念的推广；（2）用测度概念解释古典概型、几何概率及概率论公理化结构中关于概率的定义。
\end{exercise}

\begin{exercise}\label{exercise:ch01-50}
试证：概率定义~\ref{definition:ch01-概率} 中的三个要求可用下列两个要求代替：

（i）$P(A)\geqslant 0$，对一切 $A\in\mathscr{F}$；

（ii）若 $A_i\in\mathscr{F}$，$i=1,2,\cdots$，两两互不相容，且 $\sum_{i=1}^{\infty}A_i=\Omega$，则 $\sum_{i=1}^{\infty}P(A_i)=1$。
\end{exercise}

\end{document}